\section{Related Work}
\label{sec:related}

\paragraph{Distributed hash tables.} Chord \citep{stoica2001chord}, Pastry \citep{rowstron2001pastry}, and Kademlia \citep{maymounkov2002kademlia} established the structured-overlay paradigm Freenet's ring inherits. These systems organize peers by hashed identifiers in a ring or XOR-metric space and provide $O(\log n)$ lookup with carefully maintained routing tables. Their consistency model is fundamentally immutable: a put-then-get operation returns the published value. Freenet replaces the immutable cell with a mutable, contract-defined one, and the explicit finger-table maintenance with the implicit small-world construction of \cref{sec:smallworld}.

\paragraph{Content-addressed storage.} IPFS \citep{benet2014ipfs} popularized hash-based naming as a deployment substrate for web content, with a peer-to-peer block exchange (Bitswap) underneath. IPFS is also a fundamentally immutable layer, with mutability provided externally (IPNS, smart contracts) and consistency left to applications. Freenet's design choice is to put mutability and consistency inside the platform via the contract abstraction; this is a strictly more opinionated position than IPFS takes.

\paragraph{CRDTs.} The theory of Conflict-free Replicated Data Types \citep{shapiro2011crdt} is the algebraic foundation we build on. Freenet's contracts realize the state-based (CvRDT) flavor with two extensions: the lattice is application-supplied rather than chosen from a catalogue, and the summary/delta protocol generalizes the implicit ``ship the whole state'' protocol of vanilla state-based CRDTs. The Antidote \citep{akkoorath2016cure} and AntidoteDB \citep{antidote} projects represent the most direct points of comparison from the database side --- both are CRDT databases for the geo-distributed setting --- but they expose a fixed CRDT vocabulary rather than letting applications define their own.

\paragraph{Anti-entropy and delta sync.} Dynamo \citep{decandia2007dynamo} and its descendants used Merkle anti-entropy to reconcile replicas; Bayou \citep{terry1995bayou} pioneered an operation-log approach with application-defined merge procedures that is conceptually the closest historical analogue to Freenet's contracts. Delta-state CRDTs \citep{almeida2018delta} formalize a state-based protocol that ships only changes between replicas, very much in the same shape as \cref{sec:sync}; the differences are that Freenet leaves the delta type and protocol fully under contract control and embeds the result in a small-world routing layer rather than a fully-connected replica set.

\paragraph{Small-world routing.} Kleinberg's 2000 result \citep{kleinberg2000small} is the theoretical underpinning of Freenet's routing layer; the original Freenet \citep{clarke2001freenet} predated and informally anticipated some of it. The accept-only-at-terminus rule is, to our knowledge, an unusual way to construct the small-world topology --- it builds the link distribution from request traffic rather than from explicit sampling --- and we are not aware of a prior analysis of routing performance under that construction.

\paragraph{Performance-aware routing.} Adaptive routing in peer-to-peer overlays has a long history; isotonic regression specifically has been used in machine-learning settings as a calibration tool \citep{deleeuw2009isotone} but, as far as we are aware, not previously applied as the per-neighbor performance model in a structured overlay's forwarding decision.

\paragraph{Blockchain platforms.} Ethereum \citep{wood2014ethereum} and successors provide programmable consistency, but pay for it with a single linearized state machine replicated globally. Freenet relaxes the global linearization in favor of per-contract eventual consistency with application-defined merge, which makes a different class of applications natural --- the interactive, conversational ones --- at the cost of being unsuitable for applications that fundamentally require a single global order on transactions.

\paragraph{Federated systems.} Matrix \citep{matrix} is the closest non-peer-to-peer system to Freenet in application surface (group messaging, shared state, identity). It is server-based and federated; its state-resolution algorithm reconciles divergent room states by an application-defined deterministic rule, which is a federated analogue of contract-defined merge. Freenet pushes the same idea into a fully peer-to-peer setting and generalizes the merge to arbitrary commutative monoids.

\paragraph{The original Freenet.} As discussed in \cref{sec:discussion}, this platform shares its name and small-world intuition with the original Freenet \citep{clarke2001freenet} and very little else; the project's commit history records the trajectory from one to the other.
